\documentclass[11pt]{article}
\usepackage[margin=1.05in]{geometry}
\usepackage{amsmath,amssymb,booktabs}
\usepackage[colorlinks=true,linkcolor=black,urlcolor=blue]{hyperref}
\setlength{\parskip}{0.35em}
\newcommand{\code}[1]{\texttt{#1}}

\title{\vspace{-2em}Boundary conditions in a discontinuous-Galerkin solver\\
\large How the weak form becomes matrix blocks, and what the flags select}
\author{}
\date{}

\begin{document}
\maketitle
\vspace{-3em}

\section{What a boundary condition \emph{is} here}

Take a conservation law
\[
  \partial_t f \;=\; -\,\partial_x F , \qquad F = F(f,\partial_x f),
\]
and a mesh of cells $I_j=[x_{j-1/2},x_{j+1/2}]$. On each cell, multiply by a test
function $v$ and integrate by parts once:
\begin{equation}
  \int_{I_j} v\,\partial_t f
  \;=\; \int_{I_j} v'\,F
  \;-\;\Big[\,v\,\widehat F\,\Big]_{x_{j-1/2}}^{x_{j+1/2}} .
  \label{eq:weak}
\end{equation}
The solution is discontinuous across cell edges, so $F$ has two values there and
the scheme must pick one. That single number, $\widehat F$, is the
\emph{numerical flux}. It is the only way neighbouring cells communicate.

Everything about boundary conditions follows from one observation about
\eqref{eq:weak}. Put $v\equiv 1$ (it is in the space) and sum over all cells.
The volume term dies because $v'=0$. Every \emph{interior} edge appears twice ---
once as the right edge of $I_j$, once as the left edge of $I_{j+1}$ --- with
opposite signs and, crucially, the \emph{same} $\widehat F$, because the numerical
flux is single-valued. They cancel in pairs. What survives is
\begin{equation}
  \frac{d}{dt}\int f \;=\; \widehat F(a) \;-\; \widehat F(b)
  \label{eq:ledger}
\end{equation}
for a domain $[a,b]$.

This is exact algebra, not an approximation: it holds at any mesh size and any
polynomial degree. The interior of the scheme cannot create or destroy anything.
\emph{All} gain and loss enters through two numbers, the numerical fluxes at the
two ends. So in a DG code:

\begin{quote}
  A boundary condition is a rule for choosing $\widehat F$ at the two domain
  edges, where there is no neighbour to supply one.
\end{quote}

Note what this is \emph{not}. It is not a condition on $f$, and not a condition
on $\partial_x f$. It is a condition on a flux. Those coincide only in special
cases, and confusing them is the most common way to get a wrong wall.

\section{Where the matrix blocks come from}

Assembling \eqref{eq:weak} over a basis of local polynomials gives, per cell, a
small dense block from the volume integral, plus blocks from the two edge terms.
An edge term couples the cell to the trace of the solution on that edge; since
each edge has a value from either side, it produces one block on the diagonal
(this cell's own trace) and one off-diagonal (the neighbour's). The result is
block-tridiagonal. The edge blocks carry a factor $1/h$: a trace evaluation of
a basis function normalised on a cell of width $h$.

Now specialise to the first and last cell. Their outward edges have no
neighbour, so the off-diagonal partner does not exist and only one block is in
question. There are two things the assembler can do with it, and they are the
two options a flag selects:

\begin{center}
\begin{tabular}{@{}lll@{}}
\toprule
option & block & meaning \\
\midrule
\emph{free} & written, using this cell's own trace &
  $\widehat F = F(\text{interior})$ \\
\emph{fixed} & omitted &
  $\widehat F = 0$ \\
\bottomrule
\end{tabular}
\end{center}

The second line is worth pausing on. Prescribing zero flux is implemented by
\emph{not writing anything}: a zero contributes nothing to the matrix. The
boundary condition is an absence. (A nonzero prescribed flux is a known number
rather than an unknown, so it goes to the right-hand side as a source instead.)

Substituting into \eqref{eq:ledger}: a fixed edge contributes nothing to
$d\!\int\! f/dt$, and a free edge contributes whatever flux the solution happens
to carry there --- an open door, whose flow is dynamic and unconstrained.

\section{First-order terms: one flux, one choice}

For an advective term $F = c\,f$ this is the whole story. Fix the edge and the
wall is sealed; leave it free and material crosses at rate $c\,f(\text{wall})$.

The sign matters more than it first appears. If $c$ points \emph{out} of the
domain at that wall, a free edge is an outflow condition: material leaves,
$\int f$ decreases. If $c$ points \emph{in}, a free edge does not drain the
domain --- it \emph{fills} it, because the scheme evaluates the flux as though
the solution continued smoothly past the wall, and that fictitious outside
region pushes its own material in. An open inflow boundary is a source. It is
not a small error and refining the mesh does not reduce it.

\section{Second-order terms: two traces, and why the flags come in pairs}

A diffusive term $F = D\,\partial_x f$ cannot be handled the same way, because
the weak form in \eqref{eq:weak} has already used its one integration by parts
and $\partial_x f$ is still sitting there, undefined across cell edges. The
standard fix (local discontinuous Galerkin) is to split the second-order
operator into two first-order ones by naming the flux as an auxiliary unknown:
\begin{equation}
  q \;=\; D\,\partial_x f, \qquad \partial_t f \;=\; -\,\partial_x q .
  \label{eq:ldg}
\end{equation}
Each of these is now first-order, and \emph{each gets its own weak form with its
own integration by parts}:
\[
  \int_{I_j} w\,q \;=\; -\!\int_{I_j} w'\,D f \;+\;\big[\,w\,D\widehat f\,\big],
  \qquad
  \int_{I_j} v\,\partial_t f \;=\; \int_{I_j} v'\,q \;-\;\big[\,v\,\widehat q\,\big].
\]
Two integrations by parts produce \textbf{two independent numerical traces}, and
they are traces of different things: $\widehat f$, the value, appears only in the
equation defining $q$; and $\widehat q$, the flux, appears only in the equation
advancing $f$.

This is the structural fact that governs everything else. In the assembled
operator the two stages appear as a matrix product,
\[
  A \;=\; A_{\text{div}}\,A_{\text{grad}},
\]
with $q$ eliminated algebraically --- it never exists as stored data, so it
cannot be given a boundary condition of its own. Each factor owns exactly one
of the two traces:
\begin{center}
\begin{tabular}{@{}lll@{}}
\toprule
factor & owns the trace & fixing it imposes \\
\midrule
the \emph{outer} one (\code{div}) & $\widehat q$ & zero \emph{flux} at that wall \\
the \emph{inner} one (\code{grad}) & $\widehat f$ & zero \emph{value} at that wall \\
\bottomrule
\end{tabular}
\end{center}

A second-order operator admits exactly one boundary condition per wall. Hence
the rule, and its two failure modes:

\begin{center}
\begin{tabular}{@{}lll@{}}
\toprule
\code{div} & \code{grad} & result at that wall \\
\midrule
fixed & free  & \textbf{Neumann}: zero flux \\
free  & fixed & \textbf{Dirichlet}: $f=0$ \\
fixed & fixed & over-determined --- two conditions, ill-posed \\
free  & free  & under-determined --- no condition \\
\bottomrule
\end{tabular}
\end{center}

The two middle rows are real, usable boundary conditions. The two outer ones are
not boundary conditions at all: they are malformed operators, and a solver will
happily assemble either and produce nonsense --- the over-determined case by
diverging, the under-determined one by drifting.

The rule is per wall, not per term: you may fix the \code{div} at one end and
the \code{grad} at the other, giving a domain that reflects on the left and
absorbs on the right.

\paragraph{A caution about vocabulary.} Some libraries flip the flag internally
on the inner factor so that the same keyword means ``impose my condition here''
on both. That is convenient but means the two factors carry \emph{opposite}
literal settings for a well-posed pair. Check the convention of the code you
are using before trusting a flag name.

\section{Several terms: only the unsealed sum leaks}

Real equations are sums of terms, $F = F_1 + F_2 + \cdots$, each assembled
separately and each carrying its own boundary choice. Because they share test
functions, the ledger \eqref{eq:ledger} sees only the total:
\begin{equation}
  \frac{d}{dt}\int f \;=\; \sum_{k\ \text{free}} \big[\widehat F_k(a) - \widehat F_k(b)\big] .
\end{equation}
Sealed terms drop out. Two consequences, both practical:

\textbf{Sealing a subset imposes one condition on their sum.} If you seal terms
$1$ and $2$ and leave $3$ free, the wall enforces $F_1+F_2 \to 0$ --- not
$F_1\to0$ and $F_2\to0$ separately --- while $F_3$ leaks. So sealing a drag term
$Af$ together with a diffusion term $D\partial_x f$ does not give $\partial_x f=0$;
it gives
\[
  A f + D\,\partial_x f \;=\; 0
  \qquad\Longrightarrow\qquad
  \frac{\partial_x f}{f} = -\frac{A}{D},
\]
a Robin condition that nobody wrote as one. For a Fokker--Planck operator this
is exactly the equilibrium slope, which is why sealing everything reproduces the
Maxwellian at the wall without anyone imposing it.

\textbf{The audit.} To find out what wall your solver actually has: list the
terms; write each one's flux; note which are sealed. The sealed ones give your
boundary condition, as a sum. The unsealed ones, summed, are your leak rate ---
evaluated at the wall value of $f$, with sign. If that sum is not what you
intended, the wall is quietly doing something else, usually without any warning
from the code.

\subsection*{A flag acts on one term only}

This is worth making explicit, because it is the reason the previous paragraph
works. The discrete operator is assembled term by term. Each term contributes
its own bilinear form, built from its own integration by parts, and they are
added:
\begin{equation}
  a(f,v) \;=\; \sum_k a_k(f,v),
  \qquad
  a_k(f,v) \;=\; \int_\Omega v'\,F_k \;-\; \sum_{\text{edges}} \big[\,v\,\widehat F_k\,\big].
  \label{eq:perterm}
\end{equation}
The boundary choice for term $k$ selects $\widehat F_k$ at the two domain edges.
It therefore appears in exactly one summand of \eqref{eq:perterm} --- changing
the flag on term $k$ changes $a_k$ and leaves every $a_{j\ne k}$ untouched, block
for block. Nothing in the assembler ever compares one term's flag with
another's, or checks that the collection is sensible.

The terms meet only afterwards, when $v\equiv1$ collapses \eqref{eq:perterm} to
the ledger: the volume integrals vanish, and what is left is the plain sum
$\sum_k \widehat F_k$ at each end. That is the whole mechanism behind the sum
rule. Flags are private to their term; their \emph{consequences} are shared,
because every term is evaluated on the same $f$ and adds into the same balance.

\section{Mixing value and flux conditions across terms}

Sections 4 and 5 can now be combined, and doing so exposes an asymmetry between
the two kinds of condition.

\begin{itemize}\itemsep2pt
\item A \textbf{flux} condition (fixing a \code{div}) removes that term's
      contribution from the ledger. Its effect is confined to its own term, and
      several of them combine additively --- $n$ sealed terms give the single
      condition $\sum_{k\in S}F_k=0$, not $n$ separate ones.
\item A \textbf{value} condition (fixing a \code{grad}) does not remove a flux.
      It drives the shared unknown $f$ to zero at that wall. Every other term
      then evaluates \emph{its} flux on a solution that is being pinned, so a
      value condition imposed through one term is felt by all of them.
\end{itemize}

Two consequences follow, and they are the practical content of this section.

\paragraph{Advective terms may be left free at a Dirichlet wall.} Their flux is
$c\,f$, and $f\to0$ there, so the flux extinguishes itself. This is why the
recipe in \S8 does not ask you to seal them: they are harmless, not overlooked.
Note this holds for advection specifically, because its flux is proportional to
the value.

\paragraph{Diffusive flux does \emph{not} vanish at a Dirichlet wall.} It is
$D\,\partial_x f$, and $f\to0$ says nothing about the slope --- indeed the slope
is generally steepest there. That surviving flux is precisely the absorption:
it is the rate at which the wall consumes material. Sealing it would give a wall
that holds $f=0$ while admitting nothing, which is not absorption but a
contradiction, and it is the ill-posed row below.

For an equation $\partial_t f = -\partial_x(A f + D\,\partial_x f)$ with the
diffusion written as a chain, the combinations at one wall are:

\begin{center}
\begin{tabular}{@{}cccll@{}}
\toprule
adv.\ \code{div} & diff.\ \code{div} & diff.\ \code{grad} & condition at the wall & character \\
\midrule
fixed & fixed & free  & $Af + D\partial_x f = 0$ & reflecting \\
free  & free  & fixed & $f = 0$                  & absorbing \\
fixed & free  & fixed & $f = 0$                  & absorbing (sealing adv.\ is redundant) \\
free  & fixed & free  & $D\partial_x f = 0$      & leaks $Af$ --- a \emph{source} if $A$ points in \\
fixed & free  & free  & $Af = 0$                 & leaks $D\partial_x f$ \\
free  & free  & free  & none                     & under-determined \\
free  & fixed & fixed & $f=0$ \emph{and} $D\partial_x f=0$ & over-determined \\
fixed & fixed & fixed & as above                 & over-determined \\
\bottomrule
\end{tabular}
\end{center}

The accounting rule that generates this table: a second-order problem admits one
\emph{independent} condition per wall. Sealed fluxes contribute one between them
(their sum). A pinned value contributes one. Whether having both is fatal
depends on whether they are independent:

\begin{itemize}\itemsep2pt
\item $f=0$ together with $Af=0$ is \emph{redundant} --- the first implies the
      second. Harmless; row 3.
\item $f=0$ together with $D\partial_x f=0$ is \emph{independent} --- value and
      slope pinned at the same point. Over-determined; rows 7--8.
\end{itemize}

Stated generally, and not only for the two-term table: fixing a \code{grad}
anywhere gives $f=0$ at that wall. You may then additionally seal any
\emph{advective} terms --- redundant, harmless --- but sealing any
\emph{diffusive} flux, whether in that chain or in a different one, adds a
genuinely independent condition and over-determines the wall. What matters is
the kind of flux you seal, not whether it belongs to the same chain as the fixed
\code{grad}.

\subsection*{Several \code{grad} conditions at once}

These do \emph{not} accumulate the way sealed \code{div}s do, and the contrast
is the cleanest way to see what each kind of flag really is.

Sealing $n$ fluxes gives one condition, $\sum_{k\in S}F_k=0$, and \emph{which}
terms you put in $S$ changes what that condition says: seal the drag alone and
you get $f=0$; seal drag and diffusion together and you get the Robin relation
of \S5. The $F_k$ are different objects, so including more of them builds a
different sum.

Fixing $n$ \code{grad}s gives one condition, $f=0$, stated $n$ times. Every
\code{grad} flag reaches for a trace of the \emph{same} field --- the value
$\widehat f$, not a term-specific flux --- and the only value it can be assigned
is zero. A second fixed \code{grad} therefore cannot say anything the first did
not, and cannot contradict it either:
\begin{equation}
  \underbrace{\widehat f = 0}_{\text{chain }1}
  \quad\text{and}\quad
  \underbrace{\widehat f = 0}_{\text{chain }2}
  \qquad\Longleftrightarrow\qquad
  \widehat f = 0 .
\end{equation}
Two \code{grad} flags never over-determine each other. What they do change is
\emph{stiffness}: each chain penalises $f(\text{wall})\ne0$ in proportion to its
own diffusivity, so enforcement is as strong as $D_1+D_2+\dots$ rather than any
one of them. Adding a second is a way to make an absorbing wall firmer, never a
way to make it a different wall.

The asymmetry in one line: \emph{flux flags compose, value flags coincide.}

A last remark on strength. The pinning is weak, not algebraic: it is imposed to
the order of the scheme, not exactly, and if the responsible $D$ is small the
wall is only weakly absorbing. If the \code{grad} carries a coefficient $g$,
what is pinned is $g\,f=0$ --- still $f=0$ wherever $g\ne0$, but nothing at all
where $g$ vanishes. Worth checking directly: coefficients like $x^2$ or
$\sin\theta$ do vanish at a coordinate origin, and a flag placed there is
decoration.

\subsection*{Penalty terms invert the convention}

One exception, and it is a trap, because it inverts the rule everything above
rests on. For a \code{div} or \code{grad}, \code{none} \emph{writes} the
boundary block (open) and a sealing flag \emph{omits} it. For a
\code{penalty} the assembler does the reverse:

\begin{center}
\begin{tabular}{@{}lll@{}}
\toprule
term kind & \code{none} & \code{bothsides}/\code{left}/\code{right} \\
\midrule
\code{div}, \code{grad} & block written from interior trace (open) & block omitted (flux $=0$) \\
\code{penalty}          & \emph{no penalty applied} (inert)        & penalty block added (Dirichlet) \\
\bottomrule
\end{tabular}
\end{center}

The reason is that a penalty is not a flux. It penalises the jump $[f]$ across
an interface, and at a domain edge there is no neighbour, hence no jump --- so
the natural thing at a free wall is to do nothing, which is what \code{none}
selects. Giving a penalty a sealing flag instead makes it penalise $f-0$, i.e.\
enforce $f=0$: a Dirichlet penalty.

The practical consequence: \emph{do not apply a uniform ``seal everything''
rule.} A penalty swept up in it stops being inert and becomes a value
condition, which --- placed alongside the sealed diffusive fluxes it was meant
to accompany --- lands you in the over-determined row of the table above. A
zero-flux wall wants sealed \code{div}s, free \code{grad}s, and penalties left
at \code{none}.

\section{Robin conditions}

The two options in \S2 are not the only ones. Instead of omitting the boundary
block or filling it from the solution, a third choice writes a block of the
prescribed form $r\,f(\text{wall})$. In the ledger this contributes
$\mp\, r f$, a flux proportional to the local density: a partially absorbing
wall, and the discrete form of $\widehat F = r f$.

Two things follow immediately from the block structure. First, the coefficient
$r$ has the units of a flux per unit $f$ --- an effective velocity --- not of a
logarithmic derivative; supplying $A/D$ where $r$ is wanted is off by a factor
of the diffusivity. Second, since the free option writes $c\,f$ and this writes
$r\,f$ into the same slot, choosing $r=-c$ cancels the free block exactly:
\emph{Robin with $r=-c(\text{wall})$ is identically the same operator as sealing
that term}. If a code offers both, they are two spellings of one thing, and
sealing is the safer spelling because it has no coefficient to get wrong.

Stability fixes the sign: the boundary term must remove energy, not add it,
which requires $r$ to describe outflow. Get it backwards and you have built
positive feedback on the wall value.

\section{Curvilinear coordinates and the Jacobian}

If the conserved quantity is $\int m(x) f\,dx$ --- $m=x^2$ in spherical speed,
$x^2\sin\theta$ in spherical velocity --- then the equation being solved is
$m\,\partial_t f = -\partial_x F$ and the flux is measured against that measure.
The coefficient you hand a term is the \emph{mass-weighted} flux coefficient,
$m$ times the physical one. This is why a physical drag $A=1/(2x^2)$ appears in
the code as the constant $1/2$: the $x^2$ of the Jacobian has cancelled it.

The practical consequence is that fluxes often vanish at a wall for geometric
rather than physical reasons. At $x=0$ the factor $x^2$ kills every flux, so the
origin needs no condition. In a pitch-angle coordinate the $\sin\theta$ vanishes
at $\theta=0,\pi$ and the poles take care of themselves. The only walls that
require thought are the ones where the measure does not vanish.

\section{Choosing conditions for a new term}

\begin{enumerate}\itemsep2pt
\item Write each term's flux, including the Jacobian factor.
\item At each wall, decide the direction: is that flux carrying material in or
      out?
\item Decide the physical law you want there: reflecting ($\sum F=0$),
      absorbing ($f=0$), partially absorbing ($F=rf$), or open outflow.
\item Reflecting: seal every term's flux --- fix the \code{div} of each. For a
      chain, fix the \code{div} and leave the \code{grad} free.
\item Absorbing: fix the \code{grad} of the diffusive chain and leave its
      \code{div} free. Advective terms may stay free; $f\to0$ at the wall makes
      their flux vanish with it.
\item Audit before running (\S5): sum the unsealed fluxes and confirm it is the
      leak you meant.
\end{enumerate}

Two diagnostics are worth building in. Track $\int m f$ over time: a
steadily-drifting total means an unsealed flux you did not intend. And track
$\max|f|$ as well, because an ill-posed pairing can conserve mass perfectly
while diverging pointwise --- conservation alone will not catch it.

\section{Summary}

\begin{itemize}\itemsep2pt
\item A boundary condition in DG is a choice of numerical flux at the two domain
      edges, nothing more. Interior edges cancel exactly, so the two end fluxes
      are the only route by which anything is gained or lost.
\item In the matrix, the choice is which boundary block gets written: filled
      from the solution (free), omitted (zero flux), or prescribed ($rf$).
\item A first-order term has one flux and one choice. A free edge on an
      \emph{inflow} wall is a source, not a leak.
\item A second-order term is a product of two factors, and the two integrations
      by parts give two traces: the outer factor owns the flux, the inner owns
      the value. Exactly one may be fixed per wall --- flux gives Neumann, value
      gives Dirichlet, both is ill-posed, neither is under-determined.
\item With several terms, only the sum of the unsealed fluxes leaks; sealing a
      subset imposes one condition on that subset's sum, which is how a Robin
      condition can appear without anyone writing one.
\end{itemize}

\end{document}
